\chapter{Introduction and Problem Statement}
\label{ch:problem}

\section{Motivation}

Persuasion is the mechanism by which a charity turns a stranger's attention
into a donation. Every fundraising call, every solicitation script, and every
automated donation-request chatbot is, underneath, a sequence of persuasive
moves -- appeals to empathy, statements of credibility, offers of
reciprocity, statements of urgency -- chosen (consciously or not) to move a
specific person toward a specific decision. Two questions follow immediately
from that observation, and both matter well beyond charity:

\begin{enumerate}
\item \textbf{Recognition.} Given the text of a persuasive dialogue, can a
  system automatically identify \emph{which} persuasion strategies are being
  used, turn by turn, at a level of granularity fine enough to be useful?
\item \textbf{Consequence.} Once strategies can be recognized at scale, do
  some of them actually predict a better outcome than others -- and can that
  be established with enough rigor to be trusted, rather than merely
  correlated with confounds like dialogue length or persuader skill?
\end{enumerate}

Answering these questions well has direct value for nonprofit fundraising
(coaching material, script optimization, live-call feedback), for
conversational-AI systems that must persuade responsibly (donation bots,
negotiation agents, health-behavior nudging), and for computational social
science, which has long theorized about persuasion \citep{cialdini1984influence}
but rarely had a large, turn-level, multi-label-annotated corpus to test
those theories against.

\section{The Persuasion-For-Good Corpus and Its Limitation}

The \emph{Persuasion-For-Good} (P4G) corpus \citep{wang2019persuasion} is the
largest public resource of real donation-solicitation dialogues: crowd-workers
were paired as a \emph{persuader}, who is asked to convince a
\emph{persuadee} to donate part of a small task payment to the charity
\emph{Save the Children}, and a \emph{persuadee}, who is free to donate any
amount from \$0 up to the full payment. The corpus records the full dialogue
transcript and the final binary/graded donation outcome for over a thousand
conversations.

What the public release does \emph{not} provide is fine-grained, exhaustive,
turn-level labeling of \emph{which} persuasion strategy each persuader turn
uses. The original paper annotates only a small demonstrative subset with a
coarse strategy label. This is the gap the project addresses directly: without
knowing which strategies actually occur, and how densely they co-occur within
a single turn, it is not possible to build either an accurate automatic
strategy recognizer or a trustworthy strategy-to-outcome model.

\section{Problem Statement}

This project is organized around four concrete sub-problems, each addressed
by a distinct module of the system described in Chapter~\ref{ch:design}:

\begin{description}
\item[P1 -- Strategy taxonomy and annotation.] No public, exhaustive,
  multi-label persuasion-strategy annotation of the P4G corpus exists.
  A taxonomy has to be designed, and every one of the 10{,}600 persuader
  turns across 1{,}017 dialogues has to be labeled with the \emph{set} of
  strategies it uses (a single turn routinely uses two or three strategies
  at once -- treating this as single-label classification would be a
  modeling error, not a simplification).
\item[P2 -- Automatic multi-label strategy recognition.] Given a persuader
  turn (and its recent dialogue context), automatically predict which of the
  41 fine-grained strategies (grouped into 11 higher-level categories) are
  present, under severe class imbalance -- some strategies occur only a
  handful of times in the entire corpus -- and under an annotation process
  that is itself imperfect and must be measured, not assumed perfect.
\item[P3 -- Donation-outcome analysis.] Given the strategy labels for a
  dialogue (gold, to avoid contaminating the outcome model with the
  classifier's own errors), estimate which strategies and strategy
  combinations are statistically associated with donation, and by how much,
  in a way that is honest about the limits of an observational,
  non-randomized corpus of this size (1{,}017 dialogues).
\item[P4 -- Donation-\emph{intent} classification and generalization.] As a
  complementary and more directly deployable capability, classify whether a
  \emph{dialogue as a whole} ends in a donation intent (and of what type --
  outright, deferred, or conditional), and test whether an architecture built
  for this closely related task transfers to a structurally similar but
  distinct persuasive domain (buyer--seller price negotiation).
\end{description}

\section{Scope and Objectives}

\noindent The concrete, measurable objectives set for this project were:

\begin{itemize}
\item Design an 11-category / 41-strategy multi-label persuasion taxonomy
  covering the range of tactics observed in the P4G corpus, and apply it to
  every persuader turn in the first 10{,}600 turns of the corpus
  (Section~\ref{sec:annotation}).
\item Build a multi-label classifier for this taxonomy that is evaluated
  honestly against the corpus's own annotation-reliability ceiling, not
  against a naive upper bound of 1.0 (Section~\ref{sec:ceiling}).
\item Push the classifier's macro-F1 -- the metric that is dragged down by
  the rarest strategies -- as high as is actually achievable given that
  ceiling, using a disciplined, hypothesis-driven, reproducible experimental
  process rather than undirected hyperparameter search
  (Chapter~\ref{ch:evaluation}).
\item Use the resulting (gold) strategy labels to quantify, with corrected
  standard errors and multiple-testing control, how much of the variance in
  donation outcome the strategy composition of a dialogue explains, and
  compare a multi-label specification against a single-label projection of
  the same data to isolate whether multi-labeling itself adds predictive
  value (Section~\ref{sec:phase4-results}).
\item Build and evaluate a companion donation-intent classifier, and test
  whether its architecture generalizes to a different persuasive domain
  (Section~\ref{sec:module-a-results}, Section~\ref{sec:generic-results}).
\end{itemize}

The remainder of this report is organized as follows.
Chapter~\ref{ch:research} situates the project against prior work on
persuasion taxonomies, the P4G corpus, multi-label classification under class
imbalance, and outcome modeling from text. Chapter~\ref{ch:requirements}
states the functional and non-functional requirements the system was designed
against. Chapter~\ref{ch:design} describes the full system architecture,
including the annotation methodology and the classifier's seventeen-iteration
design history. Chapter~\ref{ch:evaluation} reports and interprets every
result, with confidence intervals and significance tests throughout.
Chapter~\ref{ch:management} covers project management, tooling, and ethical
considerations, and Chapter~\ref{ch:conclusion} summarizes contributions,
limitations, and future work.
